\section{Discussion}
\label{sec:discussion}

This study examines what connectivity-based classification of autism measures on ABIDE-I,
and which analysis choices determine its reported performance. To this end, a single
pipeline is used in which every learned quantity is fitted within the training partition.
Under this protocol, six findings are obtained. First, \reviewhl{discrimination varies across cognitive-ability subgroups in this cohort}. Second, somatomotor hypoconnectivity in autism strengthens,
rather than weakens, when head motion is matched within site. Third, feature selection before
the data are split inflates discrimination at the level of connections but not of regions.
Fourth, linear classifiers outperform fixed-configuration tree ensembles on identical
tangent-space features. Fifth, classifier scores are not associated with the degree of autism
within the autistic group. Sixth, the advantage of windowed connectivity is explained by the
covariance shrinkage that this description applies.

These findings matter for the interpretation of connectivity-based autism research. Because
\reviewhl{discrimination varies across cognitive-ability subgroups}, an accuracy reported on ABIDE-I partly reflects the
cognitive composition of the autistic sample. Benchmarks obtained on differently composed
samples are, therefore, not directly comparable. In addition, the somatomotor difference is
not explained by the recorded head motion, because it strengthens once motion is matched.
\reviewhl{This persistence under matching on recorded mean FD is preliminary and does not establish a biomarker.} Moreover, the
feature-selection result shows that reported accuracies \reviewhl{are interpreted together with} the
procedure that produced them. In addition, simple linear models remain a strong default for
high-dimensional connectivity features. The symptom result, in turn, shows that a score which
separates the diagnostic groups does not thereby measure the degree of autism. Finally, the
shrinkage control shows that an apparent gain from temporal information \reviewhl{is potentially produced} by
regularisation alone.

\subsection{Interpretation of the results}
\label{subsec:discussion-results}

Discrimination is maintained when a whole acquisition site is withheld from training. The
primary ensemble reaches a participant-pooled AUROC of 0.776 under pooled ten-fold validation
and 0.774 under leave-one-site-out validation (Table~\ref{tab:performance}). Because the
two values are almost identical, the model does not rely on knowledge of the acquisition
site. In addition, a random-effects mean AUROC of 0.803 is obtained across the 20 sites. The
\reviewhl{corresponding descriptive interval} extends from 0.725 to 0.882. \reviewhl{Because the site estimates share overlapping training data, this is not treated as a validated prediction interval for a new centre.} However, site-level AUROC ranges from 0.589 to 0.912, and the intervals of five sites
include 0.5. This spread is expected for sites of about thirty participants, so it is not
read as site-dependent failure.

Accuracy at the published level is reached without any step fitted before the data are
split. Specifically, a participant-pooled accuracy of 0.706 (95\% Wilson interval
0.678--0.733) is obtained in the 1,032-participant cohort. In comparison, accuracies of
0.700, 0.703, and 0.708 are reported by \citet{Heinsfeldetal2018}, \citet{Eslamietal2019},
and \citet{AlmuqhimSaeed2021}. Because these values lie within one interval of each other,
the comparison indicates parity rather than superiority.

The dependence of detection on cognitive ability is the first finding that concerns autism
itself. Against one fixed group of 468 non-autistic participants, AUROC falls from 0.873
below full-scale IQ 85 to 0.719 at 115 or above. Within the autistic group, classifier scores
also decline with full-scale IQ ($\rho=-0.18$, $p=0.0004$). In contrast, no such association
is found within the non-autistic group ($\rho=0.05$, $p=0.32$). The within-group association
is the more direct evidence, because it involves no comparison between the diagnostic
groups. When the groups are matched on IQ, however, discrimination remains highest in the
lowest quartile but no longer declines across the upper three. Part of the band gradient,
therefore, reflects the IQ difference between the groups. Nevertheless, an association
between classifier scores and verbal IQ has also been reported in ABIDE data
\citep{Nielsenetal2013}. \reviewhl{These exploratory results motivate direct interaction testing and independent validation before any IQ-specific reliability claim.}

The somatomotor finding shows why motion control changes the interpretation of group
differences. Head motion is known to alter measured functional connectivity
\citep{Poweretal2012}. In this pipeline, mean framewise displacement is positively
associated with mean connection strength ($\rho=0.51$ across participants and $0.45$ within
site). Because autistic participants move more, residual motion would, on average, raise
their connectivity, which is opposite to the observed reduction. Consistent with this
reasoning, the number of system pairs with reduced connectivity rises from 10 to 26 of 36 in
256 within-site pairs matched for motion. Similarly, the mean effect of the 21 somatomotor
regions strengthens from $-0.215$ in all 863 participants to $-0.249$ in 280 matched pairs.
Moreover, all 21 regions remain significant under this matching, and every one of them
strengthens. The direction agrees with
the somatomotor involvement reported by \citet{Chenetal2015} and with the reductions reported
by \citet{Holigaetal2019} and \citet{Ilioskaetal2023}. These studies, however, drew on
samples that include ABIDE data, so the agreement is not an independent replication.

The feature-selection result shows how reported accuracies \reviewhl{differ widely} under
apparently similar pipelines. When 2,000 of 19,900 connections are selected with all labels
before the split, fold-mean AUROC is raised from 0.756 to 0.987 on identical folds. In
addition, accuracy is raised from 0.692 to 0.948. In contrast, selecting 10 to 100 of 196
regions with all labels changes AUROC by only about 0.006 on average. The inflation,
therefore, depends on the granularity of the features. This dependence is plausible, because
selection among many candidate connections \reviewhl{risks exploiting} chance associations with the labels
of held-out participants. The two comparisons, however, also differ in cohort, shrinkage, and
ranking statistic, so they contrast two levels rather than establish a trend. For this
reason, a study that selects connections before cross-validation \reviewhl{risks reporting} inflated
estimates, and its procedure \reviewhl{is reported together with its accuracy}.

The classifier comparison indicates that linear models suit high-dimensional connectivity
features. On identical features and folds, the five linear classifiers reach
participant-pooled AUROCs of 0.735--0.757. In contrast, the three tree ensembles reach
0.611--0.624 (unadjusted DeLong $p\le2.4\times10^{-12}$ for the logistic model against each
ensemble). This pattern is consistent with the benchmark of \citet{Dadietal2019}, in which
tangent-space connectivity performs well with regularised linear classifiers. However, the
tree ensembles use fixed configurations, so the gap characterises these configurations
rather than the best attainable tree model.

The symptom analysis shows what the classifier score does not measure. Across both diagnostic
groups, 3 of 6 associations between the score and symptom severity survive correction,
whereas 0 of 6 survive within the autistic group. For the SRS total, for example, the
correlation is 0.438 across both groups but 0.050 within the autistic group. This contrast is
expected, because few non-autistic participants complete these instruments, and those who do
score at the low end of the scale. Any model that separates the groups therefore produces a
correlation across both groups largely by construction. The score is, for this reason, read
as a group-separating measurement rather than as an index of the degree of autism. This
distinction matters, because such correlations are often presented as evidence that a model
captures a graded phenotype.

The shrinkage control shows that an apparent advantage of windowed connectivity need not
arise from temporal information. An additional covariance shrinkage of 0.15 is applied by the
windowed description to its window covariances, whereas no such shrinkage is received by the
whole-scan description. When the same shrinkage is applied to the whole-scan estimate, its
AUROC is raised by 0.013 (unadjusted DeLong $p=0.004$). The windowed description then differs
from this matched baseline by $-0.002$ ($p=0.709$), against 0.010 ($p=0.281$) from the
unshrunk baseline. The small advantage is, therefore, attributed to the regularisation rather
than to variation within the scan. Comparisons between dynamic and static connectivity \reviewhl{accordingly require matched regularisation} of both descriptions.

Several secondary results bound what these data support. The ensemble probabilities have \reviewhl{low ten-bin expected calibration error}, with an expected calibration error of 0.046 under pooled ten-fold validation and
0.040 under leave-one-site-out validation. \reviewhl{This metric alone does not establish calibrated probabilities.} Selective prediction, however, does not provide a useful
guarantee on the answered cases. Specifically, the per-site abstention rule meets its target
but answers only 170 of 871 participants (19.5\%), and 143 of them are at a single site.
Finally, the attainable accuracy against these labels lies between approximately
0.85 and 0.99. However, this ceiling cannot be fixed from two instruments applied to one
population.

\subsection{Limitations}
\label{subsec:limitations}

The conclusions are limited by the composition of the cohort. Participants range from 6.5 to
58 years, with a median of 14.8 years, so the results do not bear on early identification.
In addition, the primary cohort contains 727 male and 144 female participants. Moreover, the
comparison group is typically developing rather than clinically referred, and every label is
a diagnosis obtained before the scan. The results, therefore, describe discrimination within
an aggregated research cohort rather than performance in a differential-diagnostic setting.

The evaluation is also limited to one dataset. Leave-one-site-out validation estimates
transfer among the acquisition sites represented in ABIDE-I. It does not, however, estimate
transfer to independent datasets, scanners, or preprocessing pipelines. For this reason,
external validation on an independent cohort is still required before any clinical use is
considered.

The group-difference analyses depend on the preprocessing and on their exploratory design.
All results are obtained with one preprocessing strategy without global signal regression,
one parcellation, and one nuisance model. The direction of the group differences is,
therefore, conditional on these choices. In addition, the group differences are univariate
contrasts computed without reference to the classifier. They describe how the groups differ,
but not which measurements the classifier uses. Moreover, the motion-matched and regional
analyses are exploratory, and the regional analysis is defined after inspection of the
system-level results.

Several design choices further limit the comparisons. Only the logistic model is tuned, so
the tree ensembles are represented by fixed configurations. In addition, the pairwise
$p$-values are unadjusted and are interpreted descriptively. Several approaches, moreover,
are evaluated but not adopted. Site harmonisation shows no evidence of benefit, within-scan
variability lowers accuracy, and three atlases perform no better than one. Similarly,
demographic variables add about 0.004 AUROC, and restriction to the most discriminating
regions lowers performance at every list length. These negative results are reported because
each of them bounds what a subsequent study on this public dataset \reviewhl{expects}.

\subsection{Additional analyses and remaining requirements}
\label{subsec:required-validation}

Four additional analyses bear on how the reported uncertainty and the cognitive-ability claim
should be read. Uncertainty is re-estimated with a site-clustered bootstrap, and the resulting
intervals are wider than the DeLong and Wilson intervals, particularly at the lower bound
(Section~\ref{subsec:results-performance}). The parametric intervals are therefore treated as
lower bounds on uncertainty rather than as complete accounts of it. Discrimination itself is not
placed in question, because no interval includes 0.5. Calibration is likewise assessed beyond the
binned error, and the probabilities are found to be correctly dispersed but slightly high on
average, while the binned error is shown to depend on the number of bins
(Section~\ref{subsec:results-secondary}). The expected calibration error alone would therefore
have overstated how well the probabilities are calibrated.

The dependence of detection on cognitive ability is also tested directly rather than inferred from
two separate within-group correlations. The interaction between diagnosis and cognitive ability is
negative under both protocols and survives correction across the IQ family for full-scale and
verbal IQ, but not for performance IQ (Section~\ref{subsec:results-cognitive-ability}). The claim
is accordingly stated for full-scale and verbal IQ, and it is not extended to performance IQ. The
analysed cohorts are additionally reconciled against the released records, so that every count in
the paper traces to an explicit inclusion rule.

Four requirements nevertheless remain open, and the claims of this study are bounded accordingly.
Preprocessing is not repeated within every inner cross-validation split, and the tangent reference
and feature scaler are fitted once on each outer-training partition rather than within each inner
split, so the nested design is complete at the outer level only. The training-fold region-selection
condition still uses a nuisance model fitted to the complete analysis cohort, so it is not a
strictly nested comparator. The permutation tests assume exchangeability under a site-stratified
null without residualising on the covariates that differ by diagnosis, and a residual-permutation
procedure such as Freedman--Lane has not yet been substituted. Finally, the pipeline has not been
validated on an independent cohort. For this reason the study is presented throughout as a
retrospective analysis of an aggregated research cohort, and no claim is made about clinical
screening, deployment, or performance at a new acquisition site.

\subsection{Future work}
\label{subsec:future-work}

Several extensions follow directly from these findings. The most direct extension is
external validation on an independent cohort, such as ABIDE-II. Because the available
ABIDE-II derivatives are produced by a different preprocessing pipeline, such a test requires
either the same preprocessing or an explicit analysis of transfer across pipelines.

The dependence of detection on cognitive ability also motivates further study. In particular,
future studies \reviewhl{could evaluate detection} in samples that are balanced for cognitive ability. In addition,
the dependence \reviewhl{could be tested} in cohorts with more cognitively able autistic participants,
more female participants, and younger children. Clinically referred comparison groups would,
moreover, move the evaluation closer to a differential-diagnostic setting.

The somatomotor difference \reviewhl{could, in turn, be tested} under alternative denoising strategies.
These strategies include global signal regression and finer measures of head motion. In
addition, tuned tree ensembles and other nonlinear models \reviewhl{could be compared} under the same
nested protocol, so that the classifier comparison extends beyond fixed configurations.

Finally, the feature-selection result suggests a reporting standard for connectivity-based
classification studies. Under this standard, every step fitted to data is placed inside the
cross-validation loop and is declared together with the reported performance. Such reporting
would allow published accuracies on ABIDE-I to be compared on an equal footing.
